\documentclass{article}

% upkg — Use Package
\usepackage{amsmath}
\usepackage{amssymb}
\usepackage{graphicx}
\usepackage{hyperref}
\usepackage{enumitem}
\usepackage{booktabs}
\usepackage{lipsum}

\title{SilkTex Feature Test}
\author{Test Author}
\date{\today}

\begin{document}
\maketitle
\tableofcontents
\newpage

% ── sec / sub / subs / par / subp ────────────────────────────────────────────

\section{Introduction}
\label{sec:intro}

This document exercises every snippet and SilkTex feature.
See Section~\ref{sec:math} for mathematics and Table~\ref{tab:data} for tabular data.
% section intro (end)

\subsection{Motivation}
\label{sub:motivation}

\lipsum[1]
% subsection motivation (end)

\subsubsection{Background}
\label{ssub:background}

\lipsum[2]
% subsubsection background (end)

\paragraph{A named paragraph}
\label{par:named}

\lipsum[3]
% paragraph named (end)

\subparagraph{Even finer detail}
\label{subp:detail}

\lipsum[4]
% subparagraph detail (end)

% ── Equations: eq / ueq / aeq / aneq / ali / gat / spl / cas ────────────────

\section{Mathematics}
\label{sec:math}

% eq — numbered equation
\begin{equation}
    E = mc^2
    \label{eq:einstein}
\end{equation}

% ueq — unnumbered display
\[
    \int_{-\infty}^{\infty} e^{-x^2}\,dx = \sqrt{\pi}
\]

% aeq — align*
\begin{align*}
    (a+b)^2 &= a^2 + 2ab + b^2 \\
    (a-b)^2 &= a^2 - 2ab + b^2
\end{align*}

% aneq — align (numbered)
\begin{align}
    \nabla \cdot \mathbf{E} &= \frac{\rho}{\varepsilon_0} \label{eq:gauss} \\
    \nabla \times \mathbf{B} &= \mu_0 \mathbf{J} + \mu_0\varepsilon_0 \frac{\partial \mathbf{E}}{\partial t} \label{eq:ampere}
\end{align}

% ali — aligned (with variant suffix)
\begin{aligned}
    f(x) &= x^2 + 3x - 4 \\
    f'(x) &= 2x + 3
\end{aligned}

% gat — gather (numbered)
\begin{gather}
    \sin^2\theta + \cos^2\theta = 1 \\
    \tan\theta = \frac{\sin\theta}{\cos\theta}
\end{gather}

% spl — split inside equation
\begin{equation}
    \begin{split}
        (x+y)^3 &= x^3 + 3x^2y \\
                &\quad + 3xy^2 + y^3
    \end{split}
\end{equation}

% frac — fraction
The quadratic formula is $x = \frac{-b \pm \sqrt{b^2-4ac}}{2a}$.

% cbrace — curly braces
The set $S = \left \{ x \in \mathbb{R} \mid x > 0 \right \}$ contains all positive reals.

% cas — cases
\begin{equation}
    |x| =
    \begin{cases}
        x,  &\text{ if } x \geq 0 \\
        -x, &\text{ if } x < 0
    \end{cases}
\end{equation}

% mat — matrix
\[
    A = \begin{pmatrix}
        1 & 2 & 3 \\
        4 & 5 & 6 \\
        7 & 8 & 9
    \end{pmatrix}
    \quad
    B = \begin{bmatrix}
        a & b \\
        c & d
    \end{bmatrix}
\]

% ── Lists: enum / item / desc / itd ──────────────────────────────────────────

\section{Lists}
\label{sec:lists}

% enum — enumerate
\begin{enumerate}
    \item First item
    \item Second item
    \item Third item
\end{enumerate}

% item — itemize
\begin{itemize}
    \item Bullet one
    \item Bullet two
    \item Bullet three
\end{itemize}

% desc — description
\begin{description}
    \item[Compiler] Runs \texttt{pdflatex} asynchronously
    \item[Preview] Poppler-based PDF viewer with SyncTeX
    \item[Snippets] VS Code-style JSON snippet engine
\end{description}

% itd — item with description (inside description env)
\begin{description}
    \item[Alpha] First letter of the Greek alphabet
    \item[Beta] Second letter of the Greek alphabet
\end{description}

% ── Tabular / tab ────────────────────────────────────────────────────────────

\section{Tables}
\label{sec:tables}

% tab — tabular inside table float
\begin{table}[htp]
    \centering
    \caption{Benchmark results}
    \label{tab:data}
    \begin{tabular}{lrr}
        \toprule
        Method    & Time (ms) & Memory (MB) \\
        \midrule
        Baseline  & 142       & 48          \\
        Optimised & 89        & 31          \\
        \bottomrule
    \end{tabular}
\end{table}

See Table~\ref{tab:data} for the numbers.

% ── Image / figure reference ──────────────────────────────────────────────────

\section{Figures}
\label{sec:figures}

% img — figure environment (placeholder filename — won't exist, tests parser)
\begin{figure}[htp]
    \centering
    % \includegraphics[scale=1.0]{example-image}
    \fbox{\rule{0pt}{4cm}\rule{8cm}{0pt}}
    \caption{A placeholder figure}
    \label{fig:placeholder}
\end{figure}

See Figure~\ref{fig:placeholder} above.

% ── Cross-references: figure / table / section / page ────────────────────────

\section{Cross-References}
\label{sec:refs}

Figure~\ref{fig:placeholder} appears on page~\pageref{fig:placeholder}.\\
Table~\ref{tab:data} appears in Section~\ref{sec:tables}.\\
Equation~(\ref{eq:einstein}) is in Section~\ref{sec:math}.

% ── Generic environment (begin) ───────────────────────────────────────────────

\section{Generic Environment}
\label{sec:env}

\begin{quote}
    The \texttt{begin} snippet inserts any arbitrary environment.
    This \texttt{quote} environment was created with it.
\end{quote}

% ── Bibliography (tests bibtex pipeline) ─────────────────────────────────────

\section{Conclusion}
\label{sec:conclusion}

\lipsum[5]

\begin{thebibliography}{9}
    \bibitem{knuth84}
        Donald E.\ Knuth,
        \textit{The \TeX{}book},
        Addison-Wesley, 1984.
    \bibitem{lamport94}
        Leslie Lamport,
        \textit{\LaTeX: A Document Preparation System},
        Addison-Wesley, 1994.
\end{thebibliography}

\end{document}
